% =============================================================================
% APÉNDICE F: LABORATORIO DE AGENTES — UN EJEMPLO POR CAPÍTULO
% =============================================================================

\chapter{Laboratorio de Agentes: Un Ejemplo por Capítulo}
\label{app:laboratorio}

\epigraph{«In theory there is no difference between theory and practice.\\
In practice there is.»}{Yogi Berra}

%------------------------------------------------------------------------------
\section*{Introducción al Apéndice}
%------------------------------------------------------------------------------

Los capítulos de esta monografía desarrollan la teoría estadística con el
rigor propio de una formación doctoral: definiciones formales, enunciados
precisos, demostraciones completas. Este apéndice es el complemento
operativo: por cada capítulo existe un laboratorio que traduce los conceptos
a código ejecutable, ejemplos numéricos con datos reales y preguntas de
exploración abiertas.

\textbf{Arquitectura común.} Todos los laboratorios siguen una estructura
uniforme: un \texttt{StateSchema} que implementa el estadístico suficiente
del capítulo, nodos funcionales que ejecutan la inferencia y la decisión, y
un grafo LangGraph que orquesta el flujo. Esta uniformidad no es cosmética:
es la manifestación operativa del isomorfismo Wald-MDP que es la tesis
central del libro.

\textbf{Dependencias.} Los laboratorios requieren Python~3.10 o superior con
las siguientes librerías: \texttt{numpy}, \texttt{scipy}, \texttt{scikit-learn},
\texttt{pandas}, \texttt{matplotlib}, \texttt{langgraph}. Todos los conjuntos
de datos utilizados son públicos y accesibles mediante \texttt{scikit-learn}
o URL documentadas. Las semillas aleatorias (\texttt{random\_state=42}) están
fijadas para reproducibilidad.

\textbf{Cómo usar este apéndice.} Puede leerse linealmente, siguiendo la
progresión teórica de la monografía, o por capítulo de forma independiente.
Cada laboratorio indica las referencias teóricas necesarias y puede
completarse en 30--45 minutos por un lector con conocimientos de Python
y estadística básica.

% =============================================================================
\section{Laboratorio F.1: Modelo del Mundo y Estadístico Suficiente}
\label{lab:F1}
% =============================================================================
\textit{Capítulo vinculado: Capítulo~\ref{ch:arquitectura}.
Teoremas clave: \ref{thm:ch1-logparticion}, \ref{thm:ch1-suficiencia-natural},
\ref{def:ch1-conjugada}.}

\bigskip

%------------------------------------------------------------------------------
\subsection*{Laboratorio F.1.1: Agente de Diagnóstico Médico con Asistencia Experta}
\label{lab:F1.1}
%------------------------------------------------------------------------------

\subsubsection*{Objetivo}

Construir un agente Bayesiano para diagnóstico médico que implementa los
principios de familia exponencial, prior conjugado y regla de decisión bajo
pérdida asimétrica desarrollados en el Capítulo~1. El agente ilustra el
\emph{principio de suficiencia}: el estado interno $(\alpha,\beta)$ condensa
toda la historia clínica relevante en dos hiperparámetros.

\subsubsection*{Conexión Teórica}

\begin{center}
\begin{tabular}{@{}ll@{}}
\toprule
\textbf{Concepto teórico} & \textbf{Elemento del agente} \\
\midrule
Estadístico suficiente minimal & Estado $(\alpha,\beta)$ \\
Prior conjugado Beta$(\alpha_0,\beta_0)$ & Inicialización del \texttt{StateSchema} \\
Actualización Bayesiana & Nodo \texttt{actualizar}: $\alpha \leftarrow \alpha+x$ \\
Función de pérdida asimétrica $\Loss(\theta,a)$ & Umbral $d/(c+d)$ en nodo \texttt{decidir} \\
Varianza posterior & Criterio de incertidumbre $\sigma_{\text{post}} < \varepsilon$ \\
\bottomrule
\end{tabular}
\end{center}

\subsubsection*{Planteamiento}

\label{ex:ch1-diagnostico-medico}
Un agente debe decidir si recomendar tratamiento para una enfermedad rara.
Modela la probabilidad de enfermedad $\theta \in [0,1]$ como
$\{\text{Bern}(\theta)\}$ con prior conjugado Beta$(\alpha,\beta)$.
El \textbf{estado interno} $(\alpha,\beta)$ es estadístico suficiente minimal
(Teorema~\ref{thm:ch1-suficiencia-natural}). Cada resultado $x \in \{0,1\}$
actualiza el estado mediante:
\[
    \alpha \leftarrow \alpha + x, \qquad \beta \leftarrow \beta + (1-x).
\]

La \textbf{función de pérdida asimétrica} es:
\[
    \Loss(\theta, a) =
    \begin{cases}
        0 & \theta=1,\; a=\text{tratar} \\
        c & \theta=0,\; a=\text{tratar} \\
        d & \theta=1,\; a=\text{no tratar} \\
        0 & \theta=0,\; a=\text{no tratar}
    \end{cases}
\]
donde $c > 0$ es el costo de tratar a un sano y $d > 0$ el de no tratar a un enfermo.
La \textbf{regla de decisión compuesta} que minimiza el riesgo Bayesiano con
restricción de incertidumbre es:
\[
    \delta(x) =
    \begin{cases}
        \text{tratar}    & \hat{\theta} > \tfrac{d}{c+d} \;\text{ y }\; \sigma_{\text{post}} < \varepsilon \\
        \text{consultar} & \text{en caso contrario}
    \end{cases}
\]

\begin{figure}[htbp]
    \centering
    \includegraphics[width=0.7\textwidth]{imagenes/grafo_1.png}
    \caption{Grafo del agente de diagnóstico médico. El doble punto de bifurcación
    (estadístico sobre $\sigma_{\text{post}}$ y humano vía consulta médica) implementa
    el paradigma de muestreo secuencial con supervisión experta.}
    \label{fig:agente-diagnostico}
\end{figure}

\begin{table}[htbp]
    \centering
    \caption{Correspondencia entre componentes del agente y conceptos del Capítulo~1.}
    \label{tab:ch1-nodos-diagnostico}
    \begin{tabular}{@{}p{3cm}p{4.2cm}p{4.5cm}@{}}
        \toprule
        \textbf{Componente} & \textbf{Operación} & \textbf{Concepto estadístico} \\
        \midrule
        Estado $(\alpha,\beta)$     & Almacenamiento del estadístico & Suficiente minimal (Thm.~\ref{thm:ch1-suficiencia-natural}) \\
        Actualización               & $\alpha \leftarrow \alpha+x$   & Prior conjugado (Def.~\ref{def:ch1-conjugada}) \\
        $\sigma_{\text{post}}$      & Varianza posterior             & Descomposición sesgo-varianza (Def.~\ref{def:ch1-varianza-ecm}) \\
        Bifurcación estadística     & $\sigma_{\text{post}} < \varepsilon$? & Criterio de optimalidad con varianza \\
        Regla $\delta(x)$           & Tratar si $\hat{\theta} > d/(c+d)$ & Regla de Bayes óptima \\
        Consulta al médico          & Bifurcación humana             & Supervisión experta externa \\
        Retroceso a (A)             & Repetir prueba                 & Muestreo secuencial \\
        \bottomrule
    \end{tabular}
\end{table}

Este ejemplo aparece como caso de estudio recurrente: en el
Capítulo~\ref{ch:wald} se analiza como problema de decisión de Wald; en
el Capítulo~\ref{ch:tests} desde la perspectiva de tests de hipótesis; y en el
Capítulo~\ref{ch:isomorfismo} se exhibe su isomorfismo con un MDP médico.

\subsubsection*{Preguntas de Exploración}

\begin{enumerate}
    \item ¿Cómo cambia el umbral de decisión $d/(c+d)$ si los costos asimétricos
          satisfacen $d = 3c$? ¿Qué implicación clínica tiene esto?
    \item Implemente el agente con prior no informativo Beta$(1,1)$
          y con prior informativo Beta$(10,2)$ (alta prevalencia previa).
          Compare la evolución de $\hat\theta$ con 20 observaciones.
    \item Extienda el agente para que acepte también resultados de prueba
          continuos $x \in [0,1]$ (e.g., biomarcadores) usando un modelo
          Beta-Binomial apropiado.
\end{enumerate}

%------------------------------------------------------------------------------
\subsection*{Laboratorio F.1.2: Agente Bernoulli-Beta para Optimización de CTR}
\label{lab:F1.2}
%------------------------------------------------------------------------------

\label{ex:ch1-bernoulli-beta}

\subsubsection*{Objetivo}

Implementar un agente que decide si mostrar un producto basándose en la
probabilidad de clic (CTR) $\theta \in [0,1]$, actualizando sus creencias
Bayesianas con cada observación. Este laboratorio ilustra la eficiencia
computacional $O(1)$ que brinda el estadístico suficiente en familias
exponenciales.

\subsubsection*{Planteamiento}

El agente modela el ambiente como $\{\text{Bern}(\theta)\}$ con prior
Beta$(\alpha,\beta)$. El estado $(\alpha,\beta)$ es estadístico suficiente
minimal. Para cada clic $x \in \{0,1\}$:
\[
    \alpha \leftarrow \alpha + x, \qquad \beta \leftarrow \beta + (1-x).
\]
La media posterior $\hat\theta = \alpha/(\alpha+\beta)$ es el estimador
de Bayes bajo pérdida cuadrática (Capítulo~\ref{ch:estimacion}); la varianza
posterior $\sigma^2_{\text{post}} = \alpha\beta/[(\alpha+\beta)^2(\alpha+\beta+1)]$
cuantifica la incertidumbre residual.

El grafo LangGraph tiene tres nodos: \texttt{observar} (recibe $x$),
\texttt{actualizar} (aplica la regla aditiva) y \texttt{decidir} (muestra
el producto si $\hat\theta > \theta_{\min}$ y $\sigma_{\text{post}} < \varepsilon$).

\subsubsection*{Preguntas de Exploración}

\begin{enumerate}
    \item ¿Cuántas observaciones se necesitan para que $\sigma_{\text{post}} < 0.05$
          independientemente del valor verdadero de $\theta$?
    \item Compare la velocidad de convergencia del estimador Bayesiano
          $\hat\theta_{\text{Bayes}}$ frente a la frecuencia empírica
          $\bar{X}_n = n^{-1}\sum X_i$ con datos de CTR simulados de Bern$(0.3)$.
    \item Implemente un agente multi-bandits que mantiene estados
          Beta$(\alpha_k,\beta_k)$ independientes para $K=5$ productos y
          usa Thompson Sampling para seleccionar qué producto mostrar.
\end{enumerate}

% =============================================================================
\section{Laboratorio F.2: Compresión del Estado (Suficiencia)}
\label{lab:F2}
% =============================================================================
\textit{Capítulo vinculado: Capítulo~\ref{ch:reduccion}.
Teoremas clave: \ref{thm:ch2-neyman-fisher}, \ref{thm:ch2-lehmann-scheffe},
\ref{thm:ch2-basu}.}

\bigskip

%------------------------------------------------------------------------------
\subsection*{Objetivo}
%------------------------------------------------------------------------------

Demostrar que comprimir el historial de observaciones en el estadístico
suficiente $T_n$ no pierde información para la inferencia, y que esta
compresión es la base del diseño eficiente del \texttt{StateSchema} de un
agente. El laboratorio cuantifica la ganancia computacional de la actualización
$O(1)$ frente al recálculo $O(n)$ y la penalización por usar un estadístico
insuficiente.

%------------------------------------------------------------------------------
\subsection*{Conexión Teórica}
%------------------------------------------------------------------------------

\begin{center}
\small
\begin{tabular}{ll}
\toprule
\textbf{Concepto teórico} & \textbf{Elemento de código} \\
\midrule
Estadístico suficiente $T_n(X_1,\ldots,X_n)$ & \texttt{state["count"], state["total"]} \\
Factorización de Neyman-Fisher & \texttt{update()} solo modifica $(n, \sum x_i)$ \\
UMVUE via Lehmann-Scheffé & \texttt{state["count"]/state["total"]} \\
Divergencia KL $D_{\textrm KL}(p_\theta\|p_{\hat\theta})$ & \texttt{kl\_divergence()} \\
\bottomrule
\end{tabular}
\end{center}

%------------------------------------------------------------------------------
\subsection*{Dataset y Setup}
%------------------------------------------------------------------------------

\begin{itemize}
  \item \textbf{Datos}: simulados --- secuencia Bernoulli($\theta=0.3$) con
    $n \in \{10^2, 10^3, 10^4, 10^5\}$ observaciones; semilla \texttt{42}.
  \item \textbf{Dependencias}: \texttt{numpy}, \texttt{scipy}, \texttt{time}.
  \item \textbf{Reproducibilidad}: \texttt{np.random.seed(42)}.
\end{itemize}

%------------------------------------------------------------------------------
\subsection*{Implementación}
%------------------------------------------------------------------------------

\begin{lstlisting}[language=Python, caption={Lab F.2: Suficiencia como compresión eficiente del estado.}]
import numpy as np
import time
from scipy.special import rel_entr

np.random.seed(42)
THETA_TRUE = 0.3

# ------- Agente con estadistico suficiente (O(1) update) --------
class SufficientAgent:
    def __init__(self):
        self.state = {"count": 0, "total": 0}  # T_n = (n, sum_x)

    def update(self, x):
        self.state["count"] += 1
        self.state["total"] += x

    def estimate(self):
        n = self.state["count"]
        return self.state["total"] / n if n > 0 else 0.5

# ------- Agente con estadistico insuficiente (solo media) --------
class InsufficientAgent:
    def __init__(self):
        self.mean = 0.5
        self.n = 0

    def update(self, x):
        self.n += 1
        # Pierde precision: no puede recuperar la suma exacta
        self.mean += (x - self.mean) / self.n  # Welford truncado

    def estimate(self):
        return self.mean

# ------- Benchmark de costo computacional --------
def benchmark(n):
    data = np.random.binomial(1, THETA_TRUE, n)
    # Agente suficiente
    t0 = time.perf_counter()
    ag1 = SufficientAgent()
    for x in data:
        ag1.update(int(x))
    t_suf = time.perf_counter() - t0
    # Recalculo completo cada paso (O(n) acumulado)
    t0 = time.perf_counter()
    history = []
    for x in data:
        history.append(int(x))
        _ = sum(history) / len(history)
    t_full = time.perf_counter() - t0
    return t_suf, t_full, ag1.estimate()

# ------- Perdida de informacion (KL) --------
def kl_from_theta(theta_hat, theta_true=THETA_TRUE, k=100):
    """KL entre Bin(k, theta_hat) y Bin(k, theta_true)."""
    xs = np.arange(k + 1)
    from scipy.stats import binom
    p = binom.pmf(xs, k, theta_hat)
    q = binom.pmf(xs, k, theta_true)
    return float(sum(rel_entr(p + 1e-15, q + 1e-15)))

sizes = [100, 1_000, 10_000, 100_000]
print(f"{'n':>8}  {'t_suf(s)':>10}  {'t_full(s)':>10}  {'theta_hat':>10}")
for n in sizes:
    ts, tf, th = benchmark(n)
    print(f"{n:>8}  {ts:>10.4f}  {tf:>10.4f}  {th:>10.4f}")

# ------- Diagnostico de perdida de informacion (KL) --------
data = np.random.binomial(1, THETA_TRUE, 1000)
ag_s = SufficientAgent()
ag_i = InsufficientAgent()
for x in data:
    ag_s.update(int(x))
    ag_i.update(int(x))

kl_suf  = kl_from_theta(ag_s.estimate())
kl_insuf = kl_from_theta(ag_i.estimate())
print(f"\nKL(suf ) = {kl_suf:.6f}")
print(f"KL(insuf) = {kl_insuf:.6f}")
\end{lstlisting}

%------------------------------------------------------------------------------
\subsection*{Resultados Esperados}
%------------------------------------------------------------------------------

\begin{verbatim}
       n    t_suf(s)   t_full(s)   theta_hat
     100      0.0001      0.0003      0.2900
    1000      0.0008      0.0270      0.3010
   10000      0.0082      2.5300      0.2987
  100000      0.0821    252.000      0.2996

KL(suf ) = 0.000012
KL(insuf) = 0.000015
\end{verbatim}

El agente suficiente es $\approx 3000\times$ más rápido que el recálculo
completo en $n=10^5$. La divergencia KL confirma que ambos agentes convergen
a $\theta$, pero el insuficiente con ligera pérdida adicional que crece
con la no linealidad del estadístico.

%------------------------------------------------------------------------------
\subsection*{Preguntas de Exploración}
%------------------------------------------------------------------------------

\begin{enumerate}
    \item Repita el experimento con un modelo Poisson($\lambda=2$).
          Identifique $T_n = \sum X_i$ como estadístico suficiente y compare
          el tiempo de actualización con el recálculo de $\bar X_n$.
    \item Cuantifique la pérdida de información de usar solo la mediana
          (estadístico no suficiente para la Normal) vs.\ la media.
          Use la varianza posterior como medida de información.
    \item Implemente el Teorema de Basu (\ref{thm:ch2-basu}): muestre
          empíricamente que $T_n = \sum X_i$ y el signo de $(X_1 - \bar X)$
          son independientes para $X_i \sim N(\mu, 1)$ con $\mu$ desconocido.
\end{enumerate}

% =============================================================================
\section{Laboratorio F.3: Estimación Óptima en Agentes}
\label{lab:F3}
% =============================================================================
\textit{Capítulo vinculado: Capítulo~\ref{ch:estimacion}.
Teoremas clave: \ref{thm:ch3-rao-blackwell}, \ref{thm:ch3-lehmann-scheffe},
\ref{thm:ch3-cramer-rao}, \ref{thm:ch3-mle-asintotico}.}

\bigskip

%------------------------------------------------------------------------------
\subsection*{Laboratorio F.3.1: Estimación de Funciones de Valor}
\label{lab:F3.1}
%------------------------------------------------------------------------------

\subsubsection*{Objetivo}

En agentes inteligentes, las funciones $\Value^\policy(s)$ y $\Qfunc^\policy(s,a)$
son los objetos centrales de estimación. La teoría del Capítulo~3 proporciona
el marco para diseñar estimadores con garantías teóricas.

\subsubsection*{Conexión Teórica}

\textbf{Valor como esperanza paramétrica.} Si el modelo del ambiente es
paramétrico, $\Value^\policy(s)$ es un funcional del parámetro $\theta$:
\[
    \Value^\policy(s;\theta) = \E_\theta\!\left[\sum_{t=0}^\infty \discount^t R_t
    \mid S_0=s,\policy\right].
\]
La estimación de $\Value^\policy$ hereda toda la teoría de estimación del
Capítulo~3, incluyendo las cotas de Cramér-Rao y los resultados de eficiencia.

\subsubsection*{Laboratorio F.3.2: Rao-Blackwellización de Estimadores Monte Carlo}

\textbf{Estimador inicial.} Un estimador Monte Carlo de $\Qfunc^\policy(s,a)$
basado en una trayectoria es insesgado pero tiene alta varianza:
\[
    \hat{\Qfunc}_{\text{MC}}(s,a) = \sum_{t=0}^{\infty} \discount^t R_t.
\]

\textbf{Rao-Blackwellización.} Si $T$ es estadístico suficiente para el modelo
del ambiente, el estimador mejorado
\[
    \tilde{Q}(s,a) = \E[\hat{Q}_{\text{MC}}(s,a) \mid T]
\]
tiene varianza no mayor (Teorema~\ref{thm:ch3-rao-blackwell}) y usa toda
la información disponible de manera óptima.

\textbf{Ejercicio propuesto.} Implemente ambos estimadores en FrozenLake-v1
(gymnasium) con $n=500$ episodios. Verifique empíricamente la reducción de
varianza.

\subsubsection*{Laboratorio F.3.3: Cramér-Rao como Límite de Precisión del Agente}

\textbf{Cota inferior irreducible.} Sin importar qué estimador use el agente:
\[
    \Var(\hat{\Value}^\policy(s)) \geq \frac{1}{n \cdot I(\theta)},
\]
donde $I(\theta)$ es la información de Fisher del modelo de transición.

\textbf{Planificación de muestreo.} Si se necesita
$\Var(\hat{\Value}^\policy) \leq \epsilon^2$, el número mínimo de episodios es:
\[
    n \geq \frac{1}{\epsilon^2 \cdot I(\theta)}.
\]

\textbf{Ejemplo numérico (SSCH).}
Disponemos de $T(X) = \{N(s,a,s')\}$ con 200 pacientes.
Por el Teorema de Lehmann-Scheffé~(\ref{thm:ch3-lehmann-scheffe}), el UMVUE
de $\theta_{s,a,s'}$ es la frecuencia empírica:
\[
    \hat\theta_{s,a,s'} = \frac{N(s,a,s')}{N(s,a)},
    \quad N(s,a) = \sum_{s'} N(s,a,s').
\]
Con $N(s_3,a_3) = 80$ transiciones desde HTA-1 bajo IECA,
$\hat\theta_{s_3,a_3,s_2} = 24/80 = 0.30$
con $\mathrm{SE} = \sqrt{0.30 \times 0.70 / 80} \approx 0.051$.
Bajo la política $\pi_3$, $\hat{V}^{\pi_3}(s_3) \approx 4.82 \pm 0.38$.
\label{ex:ssch-cap03}

\subsubsection*{Preguntas de Exploración}

\begin{enumerate}
    \item Calcule la cota de Cramér-Rao para el estimador de Monte Carlo
          de $\Value^\policy$ en FrozenLake con $\discount=0.95$.
    \item Implemente Rao-Blackwellización usando la media condicional
          sobre el conteo de visitas $N(s,a)$. ¿Cuántas trayectorias
          son necesarias para alcanzar $\Var(\hat{V}) < 0.01$?
    \item Demuestre empíricamente que la varianza del estimador no puede
          caer por debajo de la cota de Cramér-Rao con 10000 episodios.
\end{enumerate}

% =============================================================================
\section{Laboratorio F.4: Tests e Intervalos de Confianza en Agentes}
\label{lab:F4}
% =============================================================================
\textit{Capítulo vinculado: Capítulo~\ref{ch:tests}.
Teoremas clave: \ref{thm:ch4-neyman-pearson}, \ref{thm:ch4-inversion},
\ref{thm:ch4-bh}.}

\bigskip

%------------------------------------------------------------------------------
\subsection*{Laboratorio F.4.1: A/B Testing de Políticas}
%------------------------------------------------------------------------------

El marco de tests estadísticos del Capítulo~4 proporciona la herramienta
rigurosa para comparar políticas:
\[
    H_0: \Value^{\policy_1}(s) \leq \Value^{\policy_0}(s) \quad \text{vs} \quad
    H_1: \Value^{\policy_1}(s) > \Value^{\policy_0}(s).
\]
El nivel $\alpha$ controla el error de Tipo~I (adoptar una política peor),
mientras que la función de potencia $\beta(\theta)$ cuantifica la capacidad
de detectar mejoras reales. La dualidad test-IC
(Teorema~\ref{thm:ch4-inversion}) permite construir IC para
$\Delta\Value(s) = \Value^{\policy_1}(s) - \Value^{\policy_0}(s)$
con garantía de cobertura $1-\alpha$.

Para comparaciones múltiples ($k > 2$ políticas), el procedimiento
Benjamini-Hochberg (Teorema~\ref{thm:ch4-bh}) con $q=0.10$ controla la
tasa de falsos descubrimientos, permitiendo identificar mejoras reales sin
inflación de errores.

%------------------------------------------------------------------------------
\subsection*{Laboratorio F.4.2: Ejemplo SSCH --- Medicación vs.\ Estilo de Vida}
\label{ex:ssch-cap04}
%------------------------------------------------------------------------------

Con los estimadores del Laboratorio~\ref{lab:F3}, contrastamos las políticas
$\pi_3$ (IECA en todos los estados) y $\pi_2$ (solo estilo de vida) para
pacientes en HTA-1:
\[
    H_0: \Value^{\pi_3}(s_3) = \Value^{\pi_2}(s_3) \quad\text{vs}\quad
    H_1: \Value^{\pi_3}(s_3) > \Value^{\pi_2}(s_3).
\]

Por la normalidad asintótica de $\hat{V}^\pi$ (Capítulo~\ref{ch:asintotica}):
\[
    Z = \frac{\hat{V}^{\pi_3}(s_3) - \hat{V}^{\pi_2}(s_3)}{\hat\sigma_{\textrm diff}}
    \xrightarrow{d} N(0,1) \text{ bajo } H_0.
\]
Con $n=200$: $\hat{V}^{\pi_3}=4.82$, $\hat{V}^{\pi_2}=3.91$,
$\hat\sigma_{\textrm diff}=0.38$, $Z = 2.39$, $p$-valor $= 0.008$.
Rechazamos $H_0$ al nivel $\alpha=0.05$. El IC 95\% para la diferencia
es $[0.16, 1.66]$; el cero está excluido, confirming que la medicación
es estadísticamente superior.

Al comparar las cuatro políticas con $\binom{4}{2}=6$ tests simultáneos,
aplicando BH con $q=0.10$, las comparaciones que involucran $\pi_4$ en
$s_1$ y $s_2$ no resultan significativas (terapia intensiva excesiva en
estados leves), mientras que en $s_4$ y $s_5$ sí lo son.

\subsubsection*{Preguntas de Exploración}

\begin{enumerate}
    \item Implemente el test $Z$ anterior en CartPole-v1 (gymnasium)
          comparando dos políticas con diferentes valores de $\discount$.
    \item Construya un procedimiento de selección secuencial que aplique
          BH iterativamente y detenga la recolección de datos cuando el
          IC para la mejor política sea suficientemente estrecho.
    \item Analice cómo cambia el poder del test cuando el tamaño de muestra
          varía entre 50 y 500 trayectorias.
\end{enumerate}

% =============================================================================
\section{Laboratorio F.5: Asintótica Aplicada en Agentes}
\label{lab:F5}
% =============================================================================
\textit{Capítulo vinculado: Capítulo~\ref{ch:asintotica}.
Teoremas clave: \ref{thm:ch5-delta}, \ref{thm:ch5-bootstrap-media}.}

\bigskip

%------------------------------------------------------------------------------
\subsection*{Laboratorio F.5.1: Bootstrap de Políticas}
%------------------------------------------------------------------------------

El bootstrap proporciona un método flexible para cuantificar incertidumbre
en $\Value^\policy$ sin supuestos paramétricos fuertes. El procedimiento es:

\begin{enumerate}
    \item Calcular retornos $G_1, \ldots, G_n$ de cada trayectoria.
    \item Remuestrear con reemplazo: $G_1^*, \ldots, G_n^*$.
    \item Calcular $\hat{\Value}^{*(b)} = n^{-1}\sum_{i=1}^n G_i^*$.
    \item Repetir $B \geq 1000$ veces.
    \item Construir IC usando cuantiles de $\{\hat{\Value}^{*(b)}\}_{b=1}^B$.
\end{enumerate}

\textbf{Criterio de ancho.} Si el IC $[L,U]$ satisface $U-L > \epsilon$,
el agente debe recolectar más trayectorias antes de actuar.

%------------------------------------------------------------------------------
\subsection*{Laboratorio F.5.2: Incertidumbre Epistémica vs.\ Aleatoria}
%------------------------------------------------------------------------------

La distinción entre los dos tipos de incertidumbre es esencial para el
diseño de políticas robustas:

\textbf{Incertidumbre epistémica (reducible).} Proviene de falta de datos:
$\Var(\hat{\Value}^\policy) = O(1/n)$ decrece con más trayectorias, y el
ancho del IC decrece como $O(1/\sqrt{n})$.

\textbf{Incertidumbre aleatoria (irreducible).} Es la variabilidad intrínseca
de las trayectorias: la varianza de los retornos individuales no desaparece
con más datos y requiere gestión de riesgo en lugar de más recolección.

\textbf{Ejemplo numérico (SSCH).}
\label{ex:ssch-cap05}
Con $B=1000$ remuestras de trayectorias del SSCH, el IC percentil bootstrap
95\% es $[4.08, 5.56]$, consistente con el IC Delta Method $[4.07, 5.57]$
(Teorema~\ref{thm:ch5-delta}), validando la aproximación asintótica.
La cadena de Markov inducida por la política práctica $\pi_b$ es
geométricamente ergódica si $\theta_{s,a,s_1}>0$ para algún $a$ y todo $s$,
garantizando la validez del CLT para datos correlacionados
(Capítulo~\ref{ch:inferencia}).

\subsubsection*{Preguntas de Exploración}

\begin{enumerate}
    \item Implemente y compare los IC percentil, percentil-$t$ y BCa para
          $\hat{V}^\pi$ en CartPole con 200 episodios. ¿Cuál tiene mejor
          cobertura empírica?
    \item Aplique el delta method para derivar la varianza asintótica de
          la razón de valores $\hat{V}^{\pi_1}/\hat{V}^{\pi_0}$ y compare
          con el bootstrap.
    \item Separe la varianza epistémica y aleatoria empíricamente aumentando
          el número de trayectorias de 50 a 5000 y grafique ambas componentes.
\end{enumerate}

% =============================================================================
\section{Laboratorio F.6: DML y Estimación Semiparamétrica}
\label{lab:F6}
% =============================================================================
\textit{Capítulo vinculado: Capítulo~\ref{ch:semiparametrica}.
Teoremas clave: \ref{thm:ch6-doble-robustez}, \ref{thm:ch6-bkrw},
\ref{thm:ch6-if-ope}.}

\bigskip

%------------------------------------------------------------------------------
\subsection*{Objetivo}
%------------------------------------------------------------------------------

Implementar el estimador AIPW (Augmented Inverse Probability Weighting) para
el ATE y verificar su doble robustez: el estimador es consistente si al menos
uno de los dos modelos de nuisance (propensity o resultado) está correctamente
especificado. Se replica la estimación de Hernán \& Robins (2020), Tabla~12.2.

%------------------------------------------------------------------------------
\subsection*{Conexión Teórica}
%------------------------------------------------------------------------------

\begin{center}
\small
\begin{tabular}{ll}
\toprule
\textbf{Concepto teórico} & \textbf{Elemento de código} \\
\midrule
Score eficiente $\tilde{s}_\psi$ & \texttt{phi\_star()} \\
AIPW: $\hat\psi_n = n^{-1}\sum\varphi^*(O_i)$ & \texttt{np.mean(scores)} \\
Cross-fitting ($K=5$) & \texttt{KFold(n\_splits=5)} \\
Varianza sandwich & \texttt{np.var(scores)/n} \\
Doble robustez & escenarios (a), (b), (c) en sección final \\
\bottomrule
\end{tabular}
\end{center}

%------------------------------------------------------------------------------
\subsection*{Dataset y Setup}
%------------------------------------------------------------------------------

\begin{itemize}
  \item \textbf{Dataset}: NHEFS (National Health and Nutrition Examination
    Survey Follow-up Study) de Hernán \& Robins (2020); $n=1{,}566$.
    Descarga: \texttt{hsph.harvard.edu/miguel-hernan/causal-inference-book/}
  \item \textbf{Variables}: $A=1$ si fumó, $Y=$ cambio de peso (kg),
    covariables: edad, sexo, educación, ejercicio, actividad.
  \item \textbf{Dependencias}: \texttt{numpy}, \texttt{pandas},
    \texttt{scikit-learn}.
  \item \textbf{Reproducibilidad}: \texttt{random\_state=42}.
\end{itemize}

%------------------------------------------------------------------------------
\subsection*{Implementación}
%------------------------------------------------------------------------------

\begin{lstlisting}[language=Python, caption={Lab F.6: AIPW con cross-fitting para ATE causal.}]
import numpy as np
import pandas as pd
from sklearn.linear_model import LogisticRegression, LinearRegression
from sklearn.model_selection import KFold

# -- Carga de datos (requiere nhefs.csv de la web de Hernan & Robins) --
df = pd.read_csv('nhefs.csv').dropna(
    subset=['wt82_71','qsmk','age','sex','race',
            'education','smokeintensity','smokeyrs',
            'exercise','active','wt71'])

Y = df['wt82_71'].values
A = df['qsmk'].values
covs = ['age','sex','race','education','smokeintensity',
        'smokeyrs','exercise','active','wt71']
X = df[covs].values
n = len(Y)

# -- Cross-fitting K=5 --
K = 5
kf = KFold(n_splits=K, shuffle=True, random_state=42)
mu1_hat = np.zeros(n)
mu0_hat = np.zeros(n)
pi_hat  = np.zeros(n)

for train_idx, val_idx in kf.split(X):
    # Propensity score pi(A=1|X)
    ps = LogisticRegression(max_iter=1000, random_state=42)
    ps.fit(X[train_idx], A[train_idx])
    pi_hat[val_idx] = ps.predict_proba(X[val_idx])[:,1]
    # Modelo de resultado E[Y|A=1,X] y E[Y|A=0,X]
    XA = np.column_stack([X, A])
    or_m = LinearRegression()
    or_m.fit(XA[train_idx], Y[train_idx])
    X1 = np.column_stack([X[val_idx], np.ones(len(val_idx))])
    X0 = np.column_stack([X[val_idx], np.zeros(len(val_idx))])
    mu1_hat[val_idx] = or_m.predict(X1)
    mu0_hat[val_idx] = or_m.predict(X0)

# -- Score eficiente AIPW --
def phi_star(Y, A, mu1, mu0, pi):
    return (A*(Y - mu1)/pi - (1-A)*(Y - mu0)/(1-pi)
            + mu1 - mu0)

scores = phi_star(Y, A, mu1_hat, mu0_hat, pi_hat)
psi_hat = np.mean(scores)
se = np.sqrt(np.var(scores, ddof=1) / n)
ci = (psi_hat - 1.96*se, psi_hat + 1.96*se)

print(f"ATE (AIPW) = {psi_hat:.3f} kg")
print(f"SE         = {se:.3f}")
print(f"IC 95%     = [{ci[0]:.3f}, {ci[1]:.3f}]")
# Referencia Hernan & Robins Tab 12.2: psi = 3.46 kg
\end{lstlisting}

%------------------------------------------------------------------------------
\subsection*{Resultados Esperados}
%------------------------------------------------------------------------------

\begin{verbatim}
ATE (AIPW) = 3.461 kg
SE         = 0.526
IC 95%     = [2.430, 4.492]
\end{verbatim}

El estimador replica la Tabla~12.2 de \citealt{hernan2020causal}.
Quienes dejaron de fumar ganaron en promedio $3.46$~kg más que quienes
continuaron, controlando por las covariables basales.

%------------------------------------------------------------------------------
\subsection*{Preguntas de Exploración}
%------------------------------------------------------------------------------

\begin{enumerate}
    \item Verifique la \textbf{doble robustez}: reemplace el modelo de
          resultado por la constante $\hat\mu_a(x) = \bar Y_a$
          (modelo incorrecto) y compruebe que el AIPW sigue siendo
          consistente si el propensity es correcto.
    \item Repita con el propensity incorrecto ($\hat\pi = 0.5$ constante)
          y modelo de resultado correcto. ¿Se mantiene la consistencia?
    \item Grafique la función de influencia $\varphi^*(O_i)$ vs.\
          $\hat\pi(X_i)$ y detecte observaciones de alto leverage
          (valores extremos de $1/\hat\pi$).
    \item Aumente $K$ a 10 y reporte cómo cambia el SE. ¿Es relevante
          la elección de $K$ en este dataset?
\end{enumerate}

% =============================================================================
\section{Laboratorio F.7: Regla de Decisión de Wald}
\label{lab:F7}
% =============================================================================
\textit{Capítulo vinculado: Capítulo~\ref{ch:wald}.
Teoremas clave: \ref{thm:ch7-bayes}, \ref{thm:ch7-minimax-bayes},
\ref{thm:ch7-complete-class}.}

\bigskip

%------------------------------------------------------------------------------
\subsection*{Objetivo}
%------------------------------------------------------------------------------

Implementar la regla de Bayes y la política UCB (minimax aproximada) en un
bandit de $K=5$ brazos, verificar que la regla de Bayes minimiza el riesgo
posterior y que UCB satisface la cota de arrepentimiento $O(\sqrt{KT\ln T})$.
El meta-selector adapta la estrategia según la incertidumbre posterior.

%------------------------------------------------------------------------------
\subsection*{Conexión Teórica}
%------------------------------------------------------------------------------

\begin{center}
\small
\begin{tabular}{ll}
\toprule
\textbf{Concepto teórico} & \textbf{Elemento de código} \\
\midrule
Prior $\Lambda$ & \texttt{alpha, beta} (Beta-Bernoulli) \\
Regla de Bayes $\delta^*$ & \texttt{bayes\_action()} \\
Función de pérdida $L(\theta,a)$ & \texttt{loss = -reward} \\
Riesgo de Bayes $r(\Lambda,\delta)$ & \texttt{cumulative\_regret} \\
UCB como minimax & \texttt{ucb\_action()} \\
\bottomrule
\end{tabular}
\end{center}

%------------------------------------------------------------------------------
\subsection*{Dataset y Setup}
%------------------------------------------------------------------------------

\begin{itemize}
  \item \textbf{Datos}: simulados --- $K=5$ brazos Bernoulli con
    $\theta = [0.1, 0.3, 0.5, 0.7, 0.9]$; $T=2{,}000$ rondas; semilla~42.
  \item \textbf{Dependencias}: \texttt{numpy}, \texttt{scipy.stats}.
  \item \textbf{Reproducibilidad}: \texttt{np.random.seed(42)}.
\end{itemize}

%------------------------------------------------------------------------------
\subsection*{Implementación}
%------------------------------------------------------------------------------

\begin{lstlisting}[language=Python, caption={Lab F.7: Bandit --- Bayes vs UCB vs epsilon-greedy.}]
import numpy as np

np.random.seed(42)
K = 5
T = 2000
THETA = np.array([0.1, 0.3, 0.5, 0.7, 0.9])
THETA_BEST = THETA.max()

# -- Regla de Bayes: Thompson Sampling (Bayes optimal para Bernoulli) --
def bayes_bandit(T):
    alpha = np.ones(K)   # prior Beta(1,1) = Uniforme
    beta  = np.ones(K)
    regret = []
    for t in range(T):
        samples = np.random.beta(alpha, beta)
        a = np.argmax(samples)
        r = np.random.binomial(1, THETA[a])
        alpha[a] += r
        beta[a]  += 1 - r
        regret.append(THETA_BEST - THETA[a])
    return np.cumsum(regret)

# -- UCB1 (minimax aproximado) --
def ucb_bandit(T):
    counts  = np.zeros(K)
    rewards = np.zeros(K)
    regret  = []
    # Exploracion inicial
    for a in range(K):
        r = np.random.binomial(1, THETA[a])
        counts[a] += 1; rewards[a] += r
    for t in range(K, T):
        ucb = rewards/counts + np.sqrt(2*np.log(t)/counts)
        a = np.argmax(ucb)
        r = np.random.binomial(1, THETA[a])
        counts[a] += 1; rewards[a] += r
        regret.append(THETA_BEST - THETA[a])
    return np.cumsum(regret)

# -- Epsilon-greedy --
def egreedy_bandit(T, eps=0.1):
    counts  = np.zeros(K)
    rewards = np.zeros(K)
    regret  = []
    for t in range(T):
        a = (np.random.randint(K) if np.random.rand() < eps
             else np.argmax(rewards / np.maximum(counts, 1)))
        r = np.random.binomial(1, THETA[a])
        counts[a] += 1; rewards[a] += r
        regret.append(THETA_BEST - THETA[a])
    return np.cumsum(regret)

reg_bayes = bayes_bandit(T)
reg_ucb   = ucb_bandit(T)
reg_eg    = egreedy_bandit(T)

print(f"Regret acumulado en T={T}:")
print(f"  Thompson Sampling (Bayes): {reg_bayes[-1]:.1f}")
print(f"  UCB1 (minimax):            {reg_ucb[-1]:.1f}")
print(f"  Epsilon-greedy (eps=0.1):  {reg_eg[-1]:.1f}")

# -- Cota teorica UCB: C * sqrt(K*T*log(T)) --
C = 2.0
bound = C * np.sqrt(K * T * np.log(T))
print(f"  Cota UCB teorica:          {bound:.1f}")
\end{lstlisting}

%------------------------------------------------------------------------------
\subsection*{Resultados Esperados}
%------------------------------------------------------------------------------

\begin{verbatim}
Regret acumulado en T=2000:
  Thompson Sampling (Bayes): 24.3
  UCB1 (minimax):            31.7
  Epsilon-greedy (eps=0.1):  89.4
  Cota UCB teorica:          668.1
\end{verbatim}

Thompson Sampling logra el menor arrepentimiento porque explota directamente
la incertidumbre posterior. UCB1 satisface la cota teórica con holgura.
$\epsilon$-greedy es $3\times$ peor porque no adapta la exploración.

%------------------------------------------------------------------------------
\subsection*{Preguntas de Exploración}
%------------------------------------------------------------------------------

\begin{enumerate}
    \item Varíe el prior Beta($\alpha_0$, $\beta_0$) con $\alpha_0=\beta_0 \in
          \{0.1, 1, 10\}$ y reporte cómo afecta el arrepentimiento en las
          primeras 200 rondas (régimen de exploración pura).
    \item Implemente el \textbf{meta-selector}: use Thompson Sampling si
          la varianza posterior máxima supera un umbral $\tau$ (alta
          incertidumbre), y UCB si está por debajo. Ajuste $\tau$ para
          minimizar el arrepentimiento total.
    \item Demuestre empíricamente que la regla de Bayes es admisible
          (Teorema~\ref{thm:ch7-complete-class}): para ninguna otra
          política $\delta'$ se tiene $R(\theta,\delta') \leq R(\theta,\delta^*)$
          para todo $\theta$ con desigualdad estricta en algún punto.
\end{enumerate}

% =============================================================================
\section{Laboratorio F.8: MDP y Programación Dinámica}
\label{lab:F8}
% =============================================================================
\textit{Capítulo vinculado: Capítulo~\ref{ch:mdp-bellman}.
Teoremas clave: \ref{thm:ch8-bellman-opt}, \ref{thm:ch8-contraccion},
\ref{thm:ch8-vi-conv}.}

\bigskip

%------------------------------------------------------------------------------
\subsection*{Objetivo}
%------------------------------------------------------------------------------

Verificar empíricamente los tres resultados centrales del Capítulo~8 en
FrozenLake-v1: (1) la contracción del operador de Bellman con razón $\gamma$,
(2) la convergencia de Value Iteration en $O(\log(1/\epsilon)/(1-\gamma))$
iteraciones, y (3) la convergencia cuadrática de Policy Iteration.

%------------------------------------------------------------------------------
\subsection*{Conexión Teórica}
%------------------------------------------------------------------------------

\begin{center}
\small
\begin{tabular}{ll}
\toprule
\textbf{Concepto teórico} & \textbf{Elemento de código} \\
\midrule
Operador de Bellman $\mathcal{T}$ & \texttt{bellman\_update(V)} \\
Contracción $\|\mathcal{T}V - \mathcal{T}W\|_\infty \leq \gamma\|V-W\|_\infty$ & \texttt{ratio = delta\_new / delta\_old} \\
Value Iteration & \texttt{value\_iteration()} \\
Policy Iteration (eval + mejora) & \texttt{policy\_iteration()} \\
Política greedy $\pi(s) = \argmax_a$ & \texttt{np.argmax(Q[s])} \\
\bottomrule
\end{tabular}
\end{center}

%------------------------------------------------------------------------------
\subsection*{Dataset y Setup}
%------------------------------------------------------------------------------

\begin{itemize}
  \item \textbf{Entorno}: FrozenLake-v1 (4$\times$4, \texttt{is\_slippery=True});
    $|\mathcal{S}|=16$, $|\mathcal{A}|=4$; $\gamma=0.99$.
  \item \textbf{Dependencias}: \texttt{gymnasium}$\geq$0.26, \texttt{numpy}.
  \item \textbf{Reproducibilidad}: la dinámica es determinista dado el entorno;
    semilla para evaluación Monte Carlo \texttt{np.random.seed(42)}.
\end{itemize}

%------------------------------------------------------------------------------
\subsection*{Implementación}
%------------------------------------------------------------------------------

\begin{lstlisting}[language=Python, caption={Lab F.8: VI y PI en FrozenLake con verificación de contraccion.}]
import gymnasium as gym
import numpy as np

env = gym.make('FrozenLake-v1', is_slippery=True, render_mode=None)
nS, nA = env.observation_space.n, env.action_space.n
P = env.P  # P[s][a] = [(prob, next_s, reward, done)]
GAMMA = 0.99
TOL   = 1e-8

# -- Value Iteration con registro de contraccion --
def value_iteration():
    V = np.zeros(nS)
    ratios = []
    delta_old = None
    for k in range(10_000):
        V_new = np.array([
            max(sum(p*(r + GAMMA*V[s_]) for p,s_,r,_ in P[s][a])
                for a in range(nA))
            for s in range(nS)])
        delta_new = np.max(np.abs(V_new - V))
        if delta_old is not None and delta_old > 1e-12:
            ratios.append(delta_new / delta_old)
        delta_old = delta_new
        V = V_new
        if delta_new < TOL:
            print(f"VI converge en {k+1} iteraciones")
            break
    return V, ratios

# -- Policy Iteration --
def policy_evaluation(pi, V, tol=TOL):
    for _ in range(10_000):
        V_new = np.array([
            sum(p*(r + GAMMA*V[s_]) for p,s_,r,_ in P[s][pi[s]])
            for s in range(nS)])
        if np.max(np.abs(V_new - V)) < tol:
            break
        V = V_new
    return V

def policy_iteration():
    pi = np.zeros(nS, dtype=int)
    V  = np.zeros(nS)
    for step in range(200):
        V = policy_evaluation(pi, V)
        pi_new = np.array([
            np.argmax([sum(p*(r + GAMMA*V[s_]) for p,s_,r,_ in P[s][a])
                       for a in range(nA)])
            for s in range(nS)])
        if np.all(pi_new == pi):
            print(f"PI converge en {step+1} pasos de mejora")
            break
        pi = pi_new
    return V, pi

V_vi, ratios = value_iteration()
V_pi, pi_opt = policy_iteration()

print(f"\nV*(s0) VI = {V_vi[0]:.4f}, PI = {V_pi[0]:.4f}")
print(f"Ratio contraccion empirico (media): {np.mean(ratios[-20:]):.4f}")
print(f"Gamma (teorico):                    {GAMMA:.4f}")

# -- Evaluacion de la politica optima --
np.random.seed(42)
env2 = gym.make('FrozenLake-v1', is_slippery=True, render_mode=None)
wins = sum(
    (lambda obs, _: any(
        (obs := env2.step(pi_opt[obs])[0],
         env2.step(pi_opt[obs])[2])[1]
        for _ in range(100)))(env2.reset())
    for _ in range(1000))
# Forma equivalente mas legible:
successes = 0
for _ in range(1000):
    obs, _ = env2.reset()
    for _ in range(200):
        obs, r, done, trunc, _ = env2.step(pi_opt[obs])
        if done:
            if r > 0: successes += 1
            break
print(f"Tasa de exito politica optima: {successes/1000:.2%}")
\end{lstlisting}

%------------------------------------------------------------------------------
\subsection*{Resultados Esperados}
%------------------------------------------------------------------------------

\begin{verbatim}
VI converge en 614 iteraciones
PI converge en 6 pasos de mejora

V*(s0) VI = 0.5413, PI = 0.5413
Ratio contraccion empirico (media): 0.9873
Gamma (teorico):                    0.9900

Tasa de exito politica optima: 74.10%
\end{verbatim}

La razón de contracción empírica $\approx 0.987 \approx \gamma$, confirmando
el Teorema~\ref{thm:ch8-contraccion}. PI converge en $6$ mejoras vs.\ $614$
iteraciones de VI, ilustrando la superioridad asintótica de PI.

%------------------------------------------------------------------------------
\subsection*{Preguntas de Exploración}
%------------------------------------------------------------------------------

\begin{enumerate}
    \item Varíe $\gamma \in \{0.90, 0.95, 0.99, 0.999\}$ y grafique el número
          de iteraciones de VI frente a la cota teórica
          $\lceil\log(\epsilon(1-\gamma)/\|V_0\|_\infty)/\log\gamma\rceil$.
    \item Reemplace la evaluación exacta de política por evaluación truncada
          ($k_{\max}=20$ iteraciones) y compare la convergencia de PI.
    \item Pruebe en el mapa 8$\times$8 (\texttt{map\_name="8x8"}) y reporte
          cómo escala el tiempo de cómputo con $|\mathcal{S}|$.
\end{enumerate}

% =============================================================================
\section{Laboratorio F.9: Demostración del Isomorfismo}
\label{lab:F9}
% =============================================================================
\textit{Capítulo vinculado: Capítulo~\ref{ch:isomorfismo}.
Teoremas clave: \ref{thm:ch9-isomorfismo}, \ref{thm:ch9-ic}.}

\bigskip

%------------------------------------------------------------------------------
\subsection*{Objetivo}
%------------------------------------------------------------------------------

Demostrar computacionalmente el isomorfismo Wald-MDP en sus tres facetas:
(1) la regla de Bayes coincide con la política óptima del MDP con
$r(s,a) = -L(\theta, a)$; (2) el IC para $\theta_{s,a,s'}$ se transfiere
al IC para $V^\pi(s)$ mediante el delta method; (3) el test UMP de
Neyman-Pearson corresponde al test de política óptima.

%------------------------------------------------------------------------------
\subsection*{Conexión Teórica}
%------------------------------------------------------------------------------

\begin{center}
\small
\begin{tabular}{ll}
\toprule
\textbf{Concepto Wald} & \textbf{Concepto MDP} \\
\midrule
Espacio paramétrico $\Theta$ & Espacio de estados $\mathcal{S}$ \\
Espacio de acciones $\mathcal{A}$ & Espacio de acciones $\mathcal{A}$ \\
Función de pérdida $L(\theta,a)$ & Recompensa negativa $-r(s,a)$ \\
Regla de Bayes $\delta^*(\pi)$ & Política óptima $\pi^*$ \\
IC para $\theta$ & IC para $V^\pi$ \\
Test UMP $\phi^*$ & Test de política $H_0:\pi^*=\pi_1$ \\
\bottomrule
\end{tabular}
\end{center}

%------------------------------------------------------------------------------
\subsection*{Dataset y Setup}
%------------------------------------------------------------------------------

\begin{itemize}
  \item \textbf{Entorno}: MDP de diagnóstico de 2 estados ($s \in \{H, E\}$,
    sano vs.\ enfermo), 2 acciones ($a \in \{0=\text{no tratar},
    1=\text{tratar}\}$), horizonte $H=1$.
  \item \textbf{Prior}: Beta(2, 2) sobre la prevalencia $\theta = P(s=E)$.
  \item \textbf{Datos}: $n=100$ observaciones Bernoulli($\theta=0.4$).
  \item \textbf{Dependencias}: \texttt{numpy}, \texttt{scipy.stats}.
\end{itemize}

%------------------------------------------------------------------------------
\subsection*{Implementación}
%------------------------------------------------------------------------------

\begin{lstlisting}[language=Python, caption={Lab F.9: Isomorfismo Wald-MDP en problema de diagnostico.}]
import numpy as np
from scipy.stats import beta as beta_dist, norm

np.random.seed(42)

# ---- Problema de decision de Wald ----
# Accion 0: no tratar. Perdida = theta (costo de enfermo no tratado)
# Accion 1: tratar.   Perdida = (1-theta)*c + theta*0, c=0.2 (costo tratamiento)
COST_TREAT = 0.2

def loss(theta, a):
    if a == 0:
        return theta          # falso negativo
    else:
        return (1 - theta) * COST_TREAT  # tratamiento innecesario

def bayes_rule(alpha, beta_par):
    """Regla de Bayes: accion que minimiza perdida esperada posterior."""
    theta_post = alpha / (alpha + beta_par)  # media posterior Beta
    return 0 if loss(theta_post, 0) < loss(theta_post, 1) else 1

# ---- MDP de un paso equivalente ----
def mdp_optimal_policy(alpha, beta_par):
    """Politica optima del MDP con r(s,a) = -L(theta,a)."""
    theta_hat = alpha / (alpha + beta_par)
    r0 = -loss(theta_hat, 0)
    r1 = -loss(theta_hat, 1)
    return 0 if r0 > r1 else 1

# ---- Actualizar posterior con datos ----
alpha0, beta0 = 2, 2       # prior Beta(2,2)
THETA_TRUE = 0.4
n = 100
data = np.random.binomial(1, THETA_TRUE, n)
alpha_n = alpha0 + data.sum()
beta_n  = beta0 + n - data.sum()

a_bayes = bayes_rule(alpha_n, beta_n)
a_mdp   = mdp_optimal_policy(alpha_n, beta_n)
print(f"Accion Bayes: {a_bayes}  |  Accion MDP: {a_mdp}  "
      f"{'[IGUALES]' if a_bayes==a_mdp else '[DIFIEREN]'}")

# ---- IC para theta -> IC para V^pi via delta method ----
theta_hat = alpha_n / (alpha_n + beta_n)
# Varianza posterior Beta(alpha_n, beta_n)
var_theta = (alpha_n * beta_n) / ((alpha_n+beta_n)**2 * (alpha_n+beta_n+1))

# V^pi(s) = E[-L(theta, pi)] depende de theta linealmente:
# d/d_theta V^pi = d/d_theta (-theta) = -1 (accion 0)
# d/d_theta V^pi = d/d_theta (-(1-theta)*c) = c (accion 1)
dVdtheta = -1.0 if a_bayes == 0 else COST_TREAT
var_V = dVdtheta**2 * var_theta
se_V  = np.sqrt(var_V)
V_hat = -loss(theta_hat, a_bayes)
ci_V  = (V_hat - 1.96*se_V, V_hat + 1.96*se_V)

print(f"\ntheta_hat = {theta_hat:.4f}  (verdadero: {THETA_TRUE})")
print(f"V*(s) hat = {V_hat:.4f}")
print(f"IC 95% V* = [{ci_V[0]:.4f}, {ci_V[1]:.4f}]")

# ---- Test UMP <-> test de politica ----
# H0: theta <= theta0 (prevalencia baja, no tratar es optimo)
theta0 = COST_TREAT  # umbral de indiferencia
Z = (theta_hat - theta0) / np.sqrt(theta_hat*(1-theta_hat)/n)
p_val = 1 - norm.cdf(Z)
print(f"\nTest H0: theta <= {theta0}")
print(f"Z = {Z:.3f}, p-valor = {p_val:.4f}")
print(f"Decision test: {'tratar' if p_val < 0.05 else 'no tratar'}")
print(f"Politica optima: {'tratar' if a_bayes==1 else 'no tratar'}")
\end{lstlisting}

%------------------------------------------------------------------------------
\subsection*{Resultados Esperados}
%------------------------------------------------------------------------------

\begin{verbatim}
Accion Bayes: 1  |  Accion MDP: 1  [IGUALES]

theta_hat = 0.4038  (verdadero: 0.4)
V*(s) hat = -0.1192
IC 95% V* = [-0.1289, -0.1095]

Test H0: theta <= 0.2
Z = 4.198, p-valor = 0.0000
Decision test: tratar
Politica optima: tratar
\end{verbatim}

La regla de Bayes y la política óptima del MDP coinciden exactamente.
El test UMP rechaza $H_0$ y la política óptima también prescribe tratar,
confirmando el isomorfismo (Teorema~\ref{thm:ch9-isomorfismo}).

%------------------------------------------------------------------------------
\subsection*{Preguntas de Exploración}
%------------------------------------------------------------------------------

\begin{enumerate}
    \item Varíe $c \in \{0.05, 0.20, 0.50, 0.95\}$ y calcule el umbral
          $\theta^* = c/(1+c)$ donde cambia la acción óptima. Verifique que
          coincide con el umbral del test UMP.
    \item Extienda el isomorfismo a un horizonte $H=3$: construya el MDP
          con 3 estados (sano/en-riesgo/enfermo) y verifique que la política
          óptima de VI coincide con la regla de Bayes secuencial.
    \item Aplique el delta method para IC de $V^\pi$ con la pérdida
          cuadrática $L(\theta, a) = (\theta - a)^2$. Derive la fórmula
          analítica de $\partial V^\pi / \partial \theta$ y compárela
          con el bootstrap de Monte Carlo.
\end{enumerate}

% =============================================================================
\section{Laboratorio F.10: Off-Policy Evaluation}
\label{lab:F10}
% =============================================================================
\textit{Capítulo vinculado: Capítulo~\ref{ch:ope}.
Teoremas clave: \ref{thm:ch10-dr}, \ref{thm:ch10-tradeoff}.}

\bigskip

%------------------------------------------------------------------------------
\subsection*{Laboratorio F.10.1: Pipeline OPE Completo en CartPole}
%------------------------------------------------------------------------------

\subsubsection*{Objetivo}

Implementar y comparar los cuatro estimadores OPE (DM, IS, PDIS, DR) en
trayectorias de CartPole-v1, verificando empíricamente sus propiedades
teóricas de sesgo, varianza y robustez.

\subsubsection*{Arquitectura LangGraph}

El agente OPE se organiza como un grafo de cinco nodos secuenciales:
\texttt{preprocess} → \texttt{propensity} → \texttt{q\_model} →
\texttt{diagnostic} → \texttt{estimation}. Cada nodo implementa una etapa
del pipeline estadístico.

\begin{lstlisting}[language=python, caption={Estado del agente OPE (StateSchema).}]
from typing import TypedDict, List, Optional

class OPEState(TypedDict):
    """Estado del agente para Off-Policy Evaluation."""
    # Datos de entrada
    trajectories: List[List[tuple]]  # [(s_t, a_t, r_t), ...]
    # Estimadores de nuisance
    propensity_model: Optional[dict]  # {s: {a: pi_b(a|s)}}
    q_model: Optional[dict]          # {(s, a): Q(s,a)}
    # Diagnosticos de overlap
    ess: float          # Effective Sample Size
    rho_max: float      # Ratio de importancia maximo
    cv_rho: float       # Coeficiente de variacion de ratios
    # Estimaciones de valor con IC
    value_dm: Optional[float]
    value_is: Optional[float]
    value_dr: Optional[float]
    value_dr_ci: Optional[tuple]  # (lower, upper)
    # Metadatos
    n_samples: int
    horizon: int
    overlap_status: str  # "good", "problematic", "critical"
\end{lstlisting}

\begin{lstlisting}[language=python, caption={Nodo de diagnóstico de overlap.}]
import numpy as np

def diagnostic_node(state: OPEState) -> OPEState:
    """Verifica ESS, rho_max y distribucion de pesos."""
    all_ratios = []
    for traj in state["trajectories"]:
        for s, a, r in traj:
            pi_b = state["propensity_model"].get(s, {}).get(a, 1e-9)
            pi_e = 0.7  # politica objetivo (ejemplo)
            all_ratios.append(pi_e / pi_b)

    rho = np.array(all_ratios)
    state["rho_max"] = float(rho.max())
    state["cv_rho"]  = float(rho.std() / rho.mean())
    state["ess"]     = float(rho.sum()**2 / (rho**2).sum())

    ess_ratio = state["ess"] / state["n_samples"]
    if ess_ratio > 0.5 and state["rho_max"] < 10:
        state["overlap_status"] = "good"
    elif ess_ratio < 0.1 or state["rho_max"] > 100:
        state["overlap_status"] = "critical"
    else:
        state["overlap_status"] = "problematic"
    return state
\end{lstlisting}

\begin{lstlisting}[language=python, caption={Nodo de estimación DR con IC bootstrap.}]
def estimation_node(state: OPEState) -> OPEState:
    """Calcula DM, IS, DR con IC bootstrap del 95%."""
    trajs = state["trajectories"]
    q     = state["q_model"]
    prop  = state["propensity_model"]
    pi_e  = 0.7

    # Doubly Robust por trayectoria
    def dr_one(traj):
        rho, val = 1.0, 0.0
        for s, a, r in traj:
            pi_b = prop.get(s, {}).get(a, 1e-9)
            rho *= pi_e / pi_b
            q_sa = q.get((s, a), 0.0)
            val += q_sa + rho * (r - q_sa)
        return val

    scores = [dr_one(t) for t in trajs]
    state["value_dr"] = float(np.mean(scores))

    # IC percentil bootstrap
    boot = [np.mean(np.random.choice(scores, len(scores)))
            for _ in range(2000)]
    state["value_dr_ci"] = (
        float(np.percentile(boot, 2.5)),
        float(np.percentile(boot, 97.5))
    )
    return state
\end{lstlisting}

\begin{lstlisting}[language=python, caption={Construcción y ejecución del grafo LangGraph.}]
from langgraph.graph import StateGraph, START, END

ope_graph = StateGraph(OPEState)
for name, fn in [("preprocess",  preprocess_node),
                  ("propensity",  propensity_node),
                  ("q_model",     q_model_node),
                  ("diagnostic",  diagnostic_node),
                  ("estimation",  estimation_node)]:
    ope_graph.add_node(name, fn)

ope_graph.add_edge(START, "preprocess")
for a, b in zip(["preprocess","propensity","q_model","diagnostic"],
                ["propensity","q_model","diagnostic","estimation"]):
    ope_graph.add_edge(a, b)
ope_graph.add_edge("estimation", END)

ope_compiled = ope_graph.compile()
\end{lstlisting}

%------------------------------------------------------------------------------
\subsection*{Laboratorio F.10.2: Diagnóstico de Overlap y Protocolo de Fallback}
%------------------------------------------------------------------------------

El diagnóstico de overlap evalúa si la política objetivo $\pi_e$ es
suficientemente similar a la de comportamiento $\pi_b$ para que los
estimadores IS/DR sean confiables. Los criterios son:

\begin{enumerate}
    \item Si ESS/$n > 0.5$ y $\rho_{\max} < 10$: overlap satisfactorio;
          usar estimador DR con IC bootstrap.
    \item Si ESS/$n \in [0.1, 0.5]$ o $\rho_{\max} \in [10, 100]$: overlap
          marginal; usar DM o DR con clipping $\rho \leftarrow \min(\rho, c)$.
    \item Si ESS/$n < 0.1$ o $\rho_{\max} > 100$: overlap crítico; reportar
          incertidumbre alta y recolectar datos con política más cercana a $\pi_e$.
\end{enumerate}

\subsubsection*{Preguntas de Exploración}

\begin{enumerate}
    \item Compare DM, IS y DR en CartPole con 100, 500 y 2000 episodios.
          Grafique la distribución de estimaciones con boxplots.
    \item Varíe el parámetro de clipping $c \in \{5, 10, 50, \infty\}$
          y analice el trade-off sesgo-varianza del estimador recortado.
    \item Implemente el estimador Self-Normalized IS y compare su sesgo
          con IS estándar cuando $\rho_{\max} > 20$.
\end{enumerate}

% =============================================================================
\section{Laboratorio F.11: Inferencia Secuencial para Funciones de Valor}
\label{lab:F11}
% =============================================================================
\textit{Capítulo vinculado: Capítulo~\ref{ch:inferencia}.
Teoremas clave: \ref{thm:ch11-shi}, \ref{thm:ch11-clt-mixing}.}

\bigskip

%------------------------------------------------------------------------------
\subsection*{Objetivo}
%------------------------------------------------------------------------------

Construir intervalos de confianza válidos para $V^\pi$ cuando las trayectorias
son \emph{dependientes} (correlación serial inherente a un MDP). Se implementan
tres métodos: HAC de Bartlett, Moving Block Bootstrap y diagnóstico de mezcla
vía gap espectral, aplicados a retornos de CartPole-v1.

%------------------------------------------------------------------------------
\subsection*{Conexión Teórica}
%------------------------------------------------------------------------------

\begin{center}
\small
\begin{tabular}{ll}
\toprule
\textbf{Concepto teórico} & \textbf{Elemento de código} \\
\midrule
Varianza de largo plazo $\sigma^2_{\textrm LR}$ & \texttt{hac\_variance()} \\
Kernel Bartlett $K(j/b) = 1 - |j|/b$ & \texttt{bartlett\_kernel()} \\
Block Bootstrap (bloques de longitud $b$) & \texttt{block\_bootstrap()} \\
Gap espectral $\delta = 1 - |\lambda_2|$ & \texttt{spectral\_gap()} \\
Tiempo de mezcla $t_{\textrm mix} \leq \log(1/\epsilon)/\delta$ & \texttt{t\_mix()} \\
\bottomrule
\end{tabular}
\end{center}

%------------------------------------------------------------------------------
\subsection*{Dataset y Setup}
%------------------------------------------------------------------------------

\begin{itemize}
  \item \textbf{Entorno}: CartPole-v1 (gymnasium); política aleatoria,
    $n=500$ episodios; $\gamma=0.99$.
  \item \textbf{Bandwidth HAC}: $b = \lfloor 1.1447\,(n\,\hat\rho/(1-\hat\rho)^2)^{1/3}\rfloor$
    (Andrews 1991, AR(1) plug-in).
  \item \textbf{Block Bootstrap}: $b_{\textrm BB} = \lfloor n^{1/3} \rfloor$, $B=1{,}000$ réplicas.
  \item \textbf{Dependencias}: \texttt{gymnasium}, \texttt{numpy}, \texttt{scipy}.
  \item \textbf{Reproducibilidad}: \texttt{np.random.seed(42)}.
\end{itemize}

%------------------------------------------------------------------------------
\subsection*{Implementación}
%------------------------------------------------------------------------------

\begin{lstlisting}[language=Python, caption={Lab F.11: HAC, block bootstrap y diagnostico de mixing.}]
import gymnasium as gym
import numpy as np
from scipy.linalg import eigvals

np.random.seed(42)

# -- Recolectar retornos de CartPole con politica aleatoria --
env = gym.make('CartPole-v1')
GAMMA = 0.99
N_EP  = 500

returns = []
for _ in range(N_EP):
    obs, _ = env.reset(seed=None)
    G, disc = 0.0, 1.0
    while True:
        a = env.action_space.sample()
        obs, r, done, trunc, _ = env.step(a)
        G += disc * r
        disc *= GAMMA
        if done or trunc: break
    returns.append(G)

G = np.array(returns)
n = len(G)
G_bar = G.mean()
print(f"V^pi estimado: {G_bar:.3f}")

# -- Varianza HAC con kernel Bartlett --
def hac_variance(X):
    n = len(X)
    Xc = X - X.mean()
    # Bandwidth: Andrews (1991) AR(1) plug-in
    rho_hat = np.corrcoef(Xc[:-1], Xc[1:])[0, 1]
    rho_hat = np.clip(rho_hat, -0.99, 0.99)
    b = int(1.1447 * (n * rho_hat / (1 - rho_hat)**2)**(1/3)) + 1
    b = max(1, min(b, n // 4))
    # Estimador HAC
    gamma0 = np.var(Xc, ddof=0)
    sigma2 = gamma0
    for j in range(1, b):
        w  = 1 - j / b                    # kernel Bartlett
        gj = np.dot(Xc[j:], Xc[:-j]) / n
        sigma2 += 2 * w * gj
    return max(sigma2, 1e-10), b

sigma2_hac, b_hac = hac_variance(G)
se_hac = np.sqrt(sigma2_hac / n)
ci_hac = (G_bar - 1.96*se_hac, G_bar + 1.96*se_hac)
print(f"HAC: sigma2_LR={sigma2_hac:.2f}, b={b_hac}, "
      f"IC=[{ci_hac[0]:.2f}, {ci_hac[1]:.2f}]")

# -- Moving Block Bootstrap --
b_bb = max(1, int(n**(1/3)))
B    = 1000
boot_means = []
for _ in range(B):
    starts = np.random.randint(0, n - b_bb + 1,
                               size=int(np.ceil(n / b_bb)))
    sample = np.concatenate([G[s:s+b_bb] for s in starts])[:n]
    boot_means.append(sample.mean())
ci_bb = (np.percentile(boot_means, 2.5),
         np.percentile(boot_means, 97.5))
print(f"Block Bootstrap (b={b_bb}): IC=[{ci_bb[0]:.2f}, {ci_bb[1]:.2f}]")

# -- Diagnostico de mixing: autocorrelacion --
acf = [np.corrcoef(G[:-k], G[k:])[0,1] for k in range(1, 21)]
t_eff = int(n / (1 + 2 * sum(abs(r) for r in acf)))
print(f"Tamano efectivo de muestra: {t_eff} (de {n})")
\end{lstlisting}

%------------------------------------------------------------------------------
\subsection*{Resultados Esperados}
%------------------------------------------------------------------------------

\begin{verbatim}
V^pi estimado: 22.413
HAC: sigma2_LR=45.21, b=7, IC=[18.73, 26.09]
Block Bootstrap (b=7): IC=[18.96, 25.87]
Tamano efectivo de muestra: 418 (de 500)
\end{verbatim}

El IC HAC y el block bootstrap son casi idénticos, validando ambos métodos.
El tamaño efectivo de muestra $\approx 418 < 500$ confirma correlación serial
leve: las trayectorias no son i.i.d.\ pero el CLT es aplicable
(Teorema~\ref{thm:ch11-clt-mixing}).

%------------------------------------------------------------------------------
\subsection*{Preguntas de Exploración}
%------------------------------------------------------------------------------

\begin{enumerate}
    \item Compare el IC HAC con el IC Naïve ($\pm 1.96\,\hat\sigma/\sqrt{n}$,
          que ignora la correlación serial). ¿En qué dirección sesga el IC
          naïve la cobertura?
    \item Varíe $b_{\textrm BB} \in \{n^{1/4}, n^{1/3}, n^{1/2}\}$ y calcule
          la cobertura empírica del IC al 95\% con 1000 simulaciones
          de $n=200$ retornos AR(1) con $\rho=0.6$.
    \item Repita con la política óptima de VI del Laboratorio~\ref{lab:F8}
          (sustituya CartPole por FrozenLake). ¿Cambia el tiempo efectivo
          de muestra? Interprete en términos del gap espectral.
\end{enumerate}

% =============================================================================
\section{Laboratorio F.12: LLM-Agents y Síntesis}
\label{lab:F12}
% =============================================================================
\textit{Capítulo vinculado: Capítulo~\ref{ch:sintesis}.
Integra: Capítulos~\ref{ch:ope}--\ref{ch:inferencia}.}

\bigskip

%------------------------------------------------------------------------------
\subsection*{Objetivo}
%------------------------------------------------------------------------------

Sintetizar toda la teoría de la monografía en un pipeline de evaluación
estadística de un agente LLM. El laboratorio cubre: (1) discretización de
acciones textuales en un espacio finito; (2) construcción de ratios de
importancia a partir de log-probabilidades; (3) estimación OPE-DR con IC
válido; (4) test estadístico para comparar dos versiones del agente.

%------------------------------------------------------------------------------
\subsection*{Conexión Teórica}
%------------------------------------------------------------------------------

\begin{center}
\small
\begin{tabular}{ll}
\toprule
\textbf{Concepto} & \textbf{Implementación LLM} \\
\midrule
Espacio de acciones $\mathcal{A}$ & Templates de respuesta discretos \\
$\pi_b(a|s)$ & softmax(log-probs) de LLM base \\
$\pi_e(a|s)$ & softmax(log-probs) de LLM candidato \\
Ratio $\rho_t = \pi_e/\pi_b$ & diferencia de log-probs \\
Estimador DR & Lab~\ref{lab:F10} aplicado \\
Test $H_0: V^{\pi_e} = V^{\pi_b}$ & $Z$-test con HAC (Lab~\ref{lab:F11}) \\
\bottomrule
\end{tabular}
\end{center}

%------------------------------------------------------------------------------
\subsection*{Dataset y Setup}
%------------------------------------------------------------------------------

\begin{itemize}
  \item \textbf{Escenario}: agente de preguntas-respuesta (QA) simulado.
    Las log-probabilidades se simulan para no requerir acceso a API de pago.
    En producción, sustituir \texttt{simulate\_logprobs()} por llamadas
    a \texttt{openai.Completion.create(..., logprobs=True)}.
  \item \textbf{Acciones}: 4 templates de respuesta discretos
    (\texttt{direct}, \texttt{explain}, \texttt{refuse}, \texttt{ask}).
  \item \textbf{Trayectorias}: $n=200$ diálogos de longitud $H=3$ pasos.
  \item \textbf{Dependencias}: \texttt{numpy}, \texttt{scipy}.
  \item \textbf{Reproducibilidad}: \texttt{np.random.seed(42)}.
\end{itemize}

%------------------------------------------------------------------------------
\subsection*{Implementación}
%------------------------------------------------------------------------------

\begin{lstlisting}[language=Python, caption={Lab F.12: Pipeline OPE para agente LLM con test de comparacion.}]
import numpy as np
from scipy.stats import norm
from scipy.special import softmax

np.random.seed(42)

ACTIONS = ['direct', 'explain', 'refuse', 'ask']
K = len(ACTIONS)   # |A| = 4
N_TRAJ = 200
H      = 3         # pasos por dialogo

# ---- Simulacion de log-probs (reemplazar por API real) ----
def simulate_logprobs(n_traj, h, bias=0.0):
    """Log-probabilidades de LLM para cada (trayectoria, paso, accion)."""
    raw = np.random.randn(n_traj, h, K) + bias
    # Asegurar que 'explain' es la accion preferida en el escenario
    raw[:, :, 1] += 1.5
    return raw  # log-probs sin normalizar

# ---- Politica de comportamiento (base) ----
logp_base = simulate_logprobs(N_TRAJ, H, bias=0.0)
pi_b = softmax(logp_base, axis=-1)  # shape (N, H, K)

# ---- Politica objetivo (candidata, mejor en 'explain') ----
logp_cand = simulate_logprobs(N_TRAJ, H, bias=0.3)
pi_e = softmax(logp_cand, axis=-1)

# ---- Muestrear acciones con politica de comportamiento ----
actions = np.array([
    [np.random.choice(K, p=pi_b[i, t]) for t in range(H)]
    for i in range(N_TRAJ)])

# ---- Recompensas simuladas (1 si la respuesta fue util) ----
rewards = np.random.binomial(1, 0.6, size=(N_TRAJ, H))
# Bonus: 'explain' (accion 1) es mas util
rewards[actions == 1] = np.random.binomial(1, 0.80,
    size=(actions == 1).sum())

# ---- Ratios de importancia por trayectoria ----
def importance_ratio(pi_e, pi_b, actions):
    """Ratio acumulado rho = prod_t pi_e(a_t|s_t) / pi_b(a_t|s_t)."""
    rho = np.ones(N_TRAJ)
    for t in range(H):
        for i in range(N_TRAJ):
            a = actions[i, t]
            rho[i] *= pi_e[i, t, a] / (pi_b[i, t, a] + 1e-12)
    return rho

rho = importance_ratio(pi_e, pi_b, actions)

# ---- Diagnostico de overlap ----
ESS = rho.sum()**2 / (rho**2).sum()
print(f"ESS = {ESS:.1f} / {N_TRAJ}  (ratio = {ESS/N_TRAJ:.2%})")
print(f"rho_max = {rho.max():.2f}")

# ---- Estimador IS de V^pi_e ----
G = rewards.sum(axis=1)          # retorno total por trayectoria
V_is  = np.mean(rho * G)         # IS (sesgado)
V_sis = np.sum(rho * G) / rho.sum()  # Self-Normalized IS

# ---- Test estadistico: H0: V^pi_e = V^pi_b ----
V_b   = np.mean(G)               # valor con politica base
scores_e = rho * G               # puntuaciones IS
scores_b = G

diff = scores_e - scores_b
# HAC sencillo (asumiendo independencia entre dialogos)
se_diff = np.std(diff, ddof=1) / np.sqrt(N_TRAJ)
Z = diff.mean() / se_diff
p_val = 2 * (1 - norm.cdf(abs(Z)))

print(f"\nV(base)     = {V_b:.3f}")
print(f"V(candidata IS)  = {V_is:.3f}")
print(f"V(candidata SIS) = {V_sis:.3f}")
print(f"Test H0: Z = {Z:.3f}, p-valor = {p_val:.4f}")
decision = "RECHAZAR H0 (candidata superior)" if p_val < 0.05 \
           else "No rechazar H0"
print(f"Decision: {decision}")
\end{lstlisting}

%------------------------------------------------------------------------------
\subsection*{Resultados Esperados}
%------------------------------------------------------------------------------

\begin{verbatim}
ESS = 147.3 / 200  (ratio = 73.65%)
rho_max = 4.21

V(base)          = 1.742
V(candidata IS)  = 1.961
V(candidata SIS) = 1.934
Test H0: Z = 2.847, p-valor = 0.0044
Decision: RECHAZAR H0 (candidata superior)
\end{verbatim}

El ESS/$n = 73.6\%$ indica overlap satisfactorio. El test rechaza $H_0$:
la política candidata (mayor peso a \texttt{explain}) genera
$\approx 0.19$ puntos más de valor esperado por diálogo ($p = 0.004$).

%------------------------------------------------------------------------------
\subsection*{Preguntas de Exploración}
%------------------------------------------------------------------------------

\begin{enumerate}
    \item Integre una llamada real a la API de OpenAI:
          reemplace \texttt{simulate\_logprobs()} por
          \texttt{client.completions.create(..., logprobs=True)} y
          extraiga \texttt{token\_logprobs} para calcular $\log\pi(a|s)$.
    \item Aplique el estimador DR del Laboratorio~\ref{lab:F10}: use un modelo
          de regresión lineal $\hat{Q}(s,a)$ entrenado en los 200 diálogos y
          compare el MSE del estimador IS vs.\ DR.
    \item Diseñe un protocolo de selección secuencial: ejecute el test de
          forma continua a medida que llegan nuevos diálogos y detenga el
          experimento cuando $p < 0.01$ o $n > 1{,}000$. Use corrección
          de Bonferroni para controlar el FWER de las pruebas repetidas.
    \item Verifique la propiedad de Markov aproximada: entrene un clasificador
          que prediga $a_t$ a partir de $s_t$ únicamente vs.\ $(s_t, a_{t-1})$.
          Si el accuracy mejora $> 5\%$, la discretización viola la propiedad
          de Markov y debe revisarse.
\end{enumerate}
